\Article{8}
Gegen diesen Beweis lässt sich Folgendes einwenden.

1. Die Regel, nach der Euler schließt, dass sich aus $2m-1$ Gleichungen alle $2m-2$ Unbekannten $\alpha,\beta$ usw., $\lambda,\mu$ usw. rational bestimmen lassen, ist keineswegs allgemein, sondern lässt sehr häufig Ausnahmen zu. Betrachtet man zum Beispiel in Artikel 3 eine der Unbekannten als bekannt und versucht, die übrigen durch sie und die gegebenen Koeffizienten rational auszudrücken, so findet man leicht, dass dies unmöglich ist und keine unbekannte Größe anders als durch eine Gleichung vom Grad $m-1$ bestimmt werden kann. Obwohl sich hier unmittelbar a priori erkennen lässt, dass dies notwendig so geschehen musste, könnte man doch mit Recht bezweifeln, ob sich die Sache nicht auch im vorliegenden Fall für manche Werte von $m$ ebenso verhält, sodass sich die Unbekannten $\alpha,\beta$ usw., $\lambda,\mu$ usw. aus $u,B,C$ usw. nur durch eine Gleichung von vielleicht höherem Grad als $2m$ bestimmen lassen. Für den Fall, dass die Gleichung $X=0$ vierten Grades ist, leitet Euler rationale Werte der Koeffizienten durch $u$ und die gegebenen Koeffizienten her; dass dies aber auch bei allen höheren Gleichungen möglich ist, bedurfte gewiss einer ausführlicheren Erklärung. --- Im Übrigen scheint es der Mühe wert, jene Formeln, die $\alpha,\beta$ usw. rational durch $u,B,C$ usw. ausdrücken, tiefer und in größter Allgemeinheit zu untersuchen. Dies und andere zur Eliminationstheorie gehörende Dinge (ein keineswegs erschöpftes Thema) werde ich bei einer anderen Gelegenheit ausführlicher behandeln.

2. Selbst wenn bewiesen worden ist, dass sich unabhängig vom Grad der Gleichung $X=0$ stets Formeln finden lassen, die $\alpha,\beta$ usw., $\lambda,\mu$ usw. rational durch $u,B,C$ usw. ausdrücken, ist doch sicher, dass diese Formeln für manche bestimmten Werte der Koeffizienten $B,C$ usw. \textit{unbestimmt} werden können. Dann ist es nicht nur unmöglich, jene Unbekannten aus $u,B,C$ usw. rational zu bestimmen, sondern in manchen Fällen entsprechen einem reellen Wert von $u$ tatsächlich keine reellen Werte von $\alpha,\beta$ usw., $\lambda,\mu$ usw. Zur Bestätigung verweise ich den Leser der Kürze halber auf Eulers Abhandlung selbst, in der die Gleichung vierten Grades auf S. 236 ausführlicher behandelt wird. Jeder wird sofort sehen, dass die Formeln für die Koeffizienten $\alpha,\beta$ unbestimmt\SourceLeaf{019} werden, wenn $C=0$ ist und für $u$ der Wert $0$ angenommen wird, und dass sich ihre Werte nicht nur nicht ohne Wurzelziehen angeben lassen, sondern nicht einmal reell sind, wenn $BB-4D$ eine negative Größe ist. Zwar lässt sich in diesem Fall leicht erkennen, dass $u$ noch andere reelle Werte hat, denen reelle Werte von $\alpha,\beta$ entsprechen; dennoch könnte man befürchten, dass die Auflösung dieser Schwierigkeit (die Euler überhaupt nicht berührt hat) bei höheren Gleichungen weit größere Mühe macht. Gewiss darf diese Sache in einem genauen Beweis keineswegs stillschweigend übergangen werden.

3. Der ausgezeichnete Euler setzt stillschweigend voraus, dass die Gleichung $X=0$ insgesamt $2m$ Wurzeln hat, und setzt ihre Summe $=0$, weil das zweite Glied in $X$ fehlt. Wie ich über diese Freiheit denke (deren sich alle Autoren zu diesem Gegenstand bedienen), habe ich schon oben in Artikel 3 erklärt. Der Satz, dass die Summe aller Wurzeln einer Gleichung ihrem ersten Koeffizienten mit umgekehrtem Vorzeichen gleich ist, scheint nur auf Gleichungen anwendbar, die Wurzeln haben. Da nun gerade durch diesen Beweis gezeigt werden soll, dass die Gleichung $X=0$ tatsächlich Wurzeln hat, scheint es nicht zulässig, ihre Existenz vorauszusetzen. Ohne Zweifel werden diejenigen, die den Fehlschluss in dieser Argumentation noch nicht durchschaut haben, antworten, \textit{hier werde nicht bewiesen, dass die Gleichung $X=0$ erfüllt werden kann} (denn das bedeutet der Ausdruck, sie habe Wurzeln), \textit{sondern nur, dass sie durch Werte von $x$ erfüllt werden kann, die in der Form $a+b\sqrt{-1}$ enthalten sind; das Erstere werde dagegen als Axiom vorausgesetzt}. Da sich aber außer der reellen und der imaginären Form $a+b\sqrt{-1}$ keine anderen Formen von Größen denken lassen, scheint nicht hinreichend klar, wie sich das, was bewiesen werden soll, von dem unterscheidet, was als Axiom vorausgesetzt wird. Selbst wenn es möglich wäre, noch andere Formen von Größen zu ersinnen, etwa die Form $F,F',F''$ usw., dürfte man dennoch nicht ohne Beweis zugestehen, dass eine Gleichung durch irgendeinen Wert von $x$ erfüllt werden kann, der entweder reell oder in der Form $a+b\sqrt{-1}$ oder in der Form $F$ oder in $F'$ usw. enthalten ist. Jenes Axiom kann daher keinen anderen Sinn haben als diesen: Jede Gleichung lässt sich \textit{entweder} durch einen reellen Wert der Unbekannten erfüllen \textit{oder} durch einen imaginären Wert in der Form $a+b\sqrt{-1}$ \textit{oder} vielleicht durch einen Wert in einer anderen, bisher unbekannten Form \textit{oder} durch einen Wert, der überhaupt in keiner Form enthalten ist. Wie sich aber solche Größen, von denen man sich nicht einmal eine Vorstellung machen kann --- ein wahrer Schatten eines Schattens ---, addieren oder multiplizieren lassen, versteht man gewiss nicht mit der Klarheit, die in der Mathematik stets gefordert wird\HistoricalNote{4}.

\SourceLeaf{020}
Im Übrigen will ich die Schlussfolgerungen, die Euler aus seiner Voraussetzung gezogen hat, durch diese Einwände keineswegs verdächtig machen. Vielmehr bin ich sicher, dass sie sich durch eine weder schwierige noch von der \textsc{Eulerschen} sehr verschiedene Methode so bestätigen lassen, dass niemandem auch nur der geringste Zweifel bleiben dürfte. Ich beanstande allein die \textit{Form}, die zwar beim \textit{Finden} neuer Wahrheiten von großem Nutzen sein kann, beim \textit{Beweisen} vor der Öffentlichkeit aber keineswegs billigenswert erscheint.

4. Zum Beweis der Behauptung, das Produkt $pqr$ usw. lasse sich aus den Koeffizienten in $X$ \textit{rational} bestimmen, hat der ausgezeichnete Euler überhaupt nichts angeführt. Alles, was er hierzu bei Gleichungen \textit{vierten Grades} erläutert, ist Folgendes (wobei $a,b,c,d$ die Wurzeln der vorgelegten Gleichung $x^4+Bxx+Cx+D=0$ sind):

\lq Man wird mir ohne Zweifel einwenden, ich hätte hier vorausgesetzt, die Größe $pqr$ sei eine reelle Größe und ihr Quadrat $ppqqrr$ sei positiv; das sei noch zweifelhaft gewesen, da die Wurzeln $a,b,c,d$ imaginär seien und es daher sehr wohl geschehen könne, dass das Quadrat der aus ihnen zusammengesetzten Größe $pqr$ negativ sei. Darauf antworte ich, dass dieser Fall niemals eintreten kann; denn wie imaginär auch die Wurzeln $a,b,c,d$ sein mögen, weiß man doch, dass $a+b+c+d=0$ gelten muss: $ab+ac+ad+bc+bd+cd=B$; $abc+abd+acd+bcd=-C$\HistoricalNote{5}; $abcd=D$, wobei diese Größen $B,C,D$ reell sind. Da aber $p=a+b$, $q=a+c$, $r=a+d$ ist, lässt sich ihr Produkt \mbox{$pqr=(a+b)(a+c)(a+d)$}, \textit{wie man weiß}, durch die Größen $B,C,D$ bestimmen und wird folglich reell sein, ebenso wie wir gesehen haben, dass tatsächlich $pqr=-C$ und \mbox{$ppqqrr=CC$} ist. Ebenso wird man leicht erkennen, dass derselbe Umstand bei höheren Gleichungen eintreten muss und man mir von dieser Seite keine Einwände machen kann.\rq{}
Die Bedingung, dass sich das Produkt $pqr$ usw. \textit{rational} durch $B,C$ usw. bestimmen lässt, hat Euler nirgends hinzugefügt; doch scheint er sie stets mitverstanden zu haben, da der Beweis ohne sie keine Beweiskraft haben kann. Nun ist zwar bei Gleichungen vierten Grades richtig, dass die Entwicklung des Produkts $(a+b)(a+c)(a+d)$ den Ausdruck $aa(a+b+c+d)+abc+abd+acd+bcd=-C$ ergibt; doch scheint nicht hinreichend klar, wie sich das Produkt bei allen höheren Gleichungen rational\SourceLeaf{021} durch die Koeffizienten bestimmen lässt. Der berühmte \textsc{de Foncenex}, der dies zuerst bemerkt hat (\textit{Miscell. phil. math. soc. Taurin.}, Bd. I, S. 117), behauptet mit Recht, dass die Methode ohne einen strengen Beweis dieses Satzes ihre ganze Beweiskraft verliert. Er gesteht aber, dass ihm dieser Beweis ziemlich schwierig erscheint, und berichtet, welchen Weg er \mbox{vergeblich} versucht hat\HistoricalNote{6}. Dennoch lässt sich die Sache ohne große Schwierigkeit durch die folgende Methode erledigen (deren wesentlichen Inhalt ich hier nur andeuten kann). Obwohl bei Gleichungen vierten Grades nicht hinreichend klar ist, dass sich das Produkt $(a+b)(a+c)(a+d)$ durch die Koeffizienten $B,C,D$ bestimmen lässt, sieht man doch leicht, dass dasselbe Produkt auch $=(b+a)(b+c)(b+d)$ sowie $=(c+a)(c+b)(c+d)$ und schließlich auch $=(d+a)(d+b)(d+c)$ ist. Das Produkt $pqr$ ist daher ein Viertel der Summe $(a+b)(a+c)(a+d)+(b+a)(b+c)(b+d)+(c+a)(c+b)(c+d)+(d+a)(d+b)(d+c)$. Dass diese nach der Entwicklung ein Polynom in den Wurzeln $a,b,c,d$ sein wird, in das alle auf dieselbe Weise eingehen, lässt sich ohne Mühe a priori erkennen. Solche Funktionen lassen sich aber stets rational durch die Koeffizienten der Gleichung ausdrücken, deren Wurzeln $a,b,c,d$ sind. --- Dasselbe ist auch klar, wenn man das Produkt $pqr$ auf die Form
\[
\tfrac12(a+b-c-d)\times\tfrac12(a+c-b-d)\times\tfrac12(a+d-b-c)
\]
bringt; denn man kann leicht voraussehen, dass dieses Produkt nach der Entwicklung alle $a,b,c,d$ auf dieselbe Weise enthalten wird. Zugleich werden Sachkundige hieraus leicht entnehmen, wie dies auf höhere Gleichungen anzuwenden ist. --- Eine vollständige Darstellung des Beweises, die ich der Kürze wegen hier nicht beifügen kann, behalte ich mir zusammen mit einer ausführlicheren Untersuchung von Funktionen, die mehrere Veränderliche auf dieselbe Weise enthalten, für eine andere Gelegenheit vor.

Im Übrigen bemerke ich, dass sich außer diesen vier Einwänden noch manches andere an Eulers Beweis beanstanden ließe. Das übergehe ich jedoch stillschweigend, um nicht etwa als allzu strenger Kritiker zu erscheinen, zumal das Vorhergehende hinreichend zu zeigen scheint, dass der Beweis jedenfalls in der von Euler vorgelegten Form keineswegs als vollständig gelten kann.

Nach diesem Beweis hat Euler noch einen anderen Weg gezeigt, den Satz für Gleichungen, deren Grad keine Zweierpotenz ist, auf die Auflösung solcher Gleichungen zurückzuführen. Da diese Methode aber für Gleichungen, deren
\SourceLeaf{022}
Grad eine Zweierpotenz ist, nichts lehrt und außerdem allen vorigen Einwänden (außer dem vierten) ebenso ausgesetzt ist wie der erste allgemeine Beweis, ist es nicht nötig, sie hier ausführlicher zu erläutern.
